\subsection{Stocks, flux et correspondance physique--économie}
\label{sec:physique}
Le parallèle entre physique et économie est le centre de la démarche, non une
analogie ajoutée après calcul. Il faut néanmoins distinguer une correspondance
de dimensions, une hypothèse de représentation et une validation empirique.
Le modèle est une étude fondamentale d'interactions : l'entité reste générique,
et peut éclairer un système concret si les mécanismes pertinents y sont
effectivement présents, sans devenir implicitement une personne puis une firme.

\begin{center}
\begin{tabularx}{\linewidth}{@{}lXX@{}}
\toprule
Objet & Dimension et opération & Lecture économique proposée \\
\midrule
$K_i$ & Stock, en joules du modèle & Taille productive mobilisable \\
$A_iK_i^{\gamma_i}$ & Production ajoutée pendant un pas, en joules & Flux productif, pas un salaire \\
$A_iK_i^{\gamma_i}/\Delta t$ & Puissance moyenne sur un pas & Capacité de production par unité de temps \\
$q$ & Principal nominal, même unité de compte que $K$ & Créance et dette, non amorties \\
$rq$ & Transfert pendant un pas & Service d'intérêt, pas création de ressource \\
$NW_i$ & Stock comptable $K_i+C_i-D_i$ & Position de bilan, distincte du stock physique \\
\bottomrule
\end{tabularx}
\end{center}
Un watt vaut un joule par seconde. Les sorties du moteur sont en joules
\emph{par pas} ; les appeler des watts nécessiterait de fixer $\Delta t$ en
secondes, ce qui n'a pas été calibré.\footnote{Dans l'écriture
$K\leftarrow K+A K^\gamma$, $A$ a pour dimension
$\mathrm{J}^{1-\gamma}$ pour la production d'un pas. Le coefficient d'une
équation différentielle correspondante aurait la dimension
$\mathrm{J}^{1-\gamma}/\mathrm{s}$. Les comparer sans préciser le pas ferait
disparaître une dimension physique. Le taux $r$ est une fraction due par pas ;
$r/\Delta t$ a la dimension inverse d'un temps.}

L'hypothèse interprétative est que certains transferts monétaires représentent
des transferts de droits à mobiliser de l'exergie. Le stock réel et le réseau
nominal sont donc couplés, mais ne sont pas identiques : un prêt déplace $K$ et
inscrit simultanément une créance et une dette. La production représente un
apport au stock modélisé ; le réservoir extérieur qui le permet n'est pas
décrit. Le bilan interne est fermé comptablement, pas thermodynamiquement sur
l'ensemble économie--environnement. En outre, l'observable dite « revenu total »
est $\mathrm{prod}+\mathrm{int\_in}$ : elle additionne production et transferts
reçus et ne doit pas être assimilée sans correction à une production nette
agrégée.

\subsection{Échelles : quels paramètres sont réellement indépendants ?}
\label{sec:parametres_reduits}
En régime homogène, sans bruit ni crédit, l'ordre production puis
dépréciation donne
\[
 K_{t+1}=(K_t+A K_t^\gamma)(1-\delta),\qquad
 K_{\rm aut}=\left[\frac{A(1-\delta)}{\delta}\right]^{1/(1-\gamma)}
\]
pour $A>0$, $0<\delta<1$ et $0<\gamma<1$. Ce point fixe déterministe n'est
pas une moyenne du système bruité. Poser $x_i=K_i/K_{\rm aut}$ transforme
la production en
\[
 x_i\longleftarrow x_i+\frac{\delta}{1-\delta}x_i^\gamma .
\]
À institution fixée, l'échelle commune de $A$ et $K_0$ peut ainsi être remplacée
par le rapport $\kappa_0=K_0/K_{\rm aut}$. Pour un toy-model, éviter de compter
deux fois la même liberté d'échelle est essentiel.\footnote{Une possibilité
équivalente est de choisir $K_0$ comme unité et de retenir
$a=A K_0^{\gamma-1}$. Il s'agit ici d'éliminer \emph{une} liberté d'unité de
capital, pas de démontrer que les autres paramètres sont redondants :
$\gamma$, $\delta$, $\sigma$, $\lambda$ et l'institution restent pertinents.
Une invariance de certaines statistiques vis-à-vis de $\lambda$ n'est pas une
équivalence trajectorielle de populations de tailles différentes.}

Le rééchelonnement $K,q,K_0\mapsto c(K,q,K_0)$ et
$A\mapsto c^{1-\gamma}A$ conserve les équations homogènes et les taux
marginaux. Les seuils numériques absolus et les arrondis du moteur limitent
l'exactitude logicielle ; le recalage du pas de temps n'est pas une symétrie
exacte du marché discret. Une hausse de $A$ à $K_0$ fixe change donc
$\kappa_0$, tandis que sa compensation $K_0\mapsto
K_0\,m^{1/(1-\gamma)}$ le préserve. Les deux sont des expériences distinctes :
la compensation est un contrôle, pas une correction physique imposée.

La tension $T=K_{\rm aut}/K_{\rm eq}$, avec
$K_{\rm eq}=(Y/(N_{\rm prod}A))^{1/\gamma}$, compare l'échelle individuelle
sans crédit à l'échelle productive observée. Elle est utile notamment à
$\delta$ fixé ; elle n'est ni une variable d'état universelle, ni un autre nom
de toute la chaîne crédit--mortalité--population.
